% 71-06(71쪽) — y=f(x) 의 그래프(극소·극대·극소가 차례로 있는 곡선)와 x축 위의 a, b, c. 스캔 화질이 낮아 TikZ 로 다시 그림.
% 원문에 있는 것만: f(a)=f(b) 를 보이는 가로 점선, a·b·c 의 세로 점선, 라벨 a, O, b, c, x, y, y=f(x). 개수를 세는 문제라 원문에 없는 값은 적지 않는다.
\begin{tikzpicture}[x=0.52cm, y=0.52cm, line width=0.5pt, font=\small]
  \draw[->, >=stealth, line width=0.7pt] (-1.3,0) -- (4.65,0) node[below=1pt] {$x$};
  \draw[->, >=stealth, line width=0.7pt] (0,-1.3) -- (0,3.6) node[left=1pt] {$y$};
  \draw[densely dashed, line width=0.4pt] (-0.35,0.95) -- (1.62,0.95);
  \draw[dotted, line width=0.4pt] (-0.35,0.95) -- (-0.35,0);
  \draw[dotted, line width=0.4pt] (1.62,0.95) -- (1.62,0);
  \draw[dotted, line width=0.4pt] (3.7,1.95) -- (3.7,0);
  \draw[line width=0.6pt, smooth] plot coordinates {(-0.95,3.1) (-0.72,2.2) (-0.52,1.5) (-0.35,0.95) (-0.12,0.48) (0.15,0.34) (0.45,0.58) (0.8,1.45) (1.15,1.76) (1.42,1.48) (1.62,0.95) (1.9,0.4) (2.15,0) (2.45,-0.38) (2.72,-0.5) (3.0,-0.28) (3.2,0.05) (3.45,0.85) (3.7,1.95) (3.95,3.2)};
  \node[below=1pt] at (-0.35,0) {$a$};
  \node[below right=-2pt and -4pt] at (0,0) {$\pt{O}$};
  \node[below=1pt] at (1.62,0) {$b$};
  \node[below=1pt] at (3.7,0) {$c$};
  \node[anchor=west] at (2.4,3.3) {$y=f(x)$};
\end{tikzpicture}
